\documentclass[10pt,a4paper]{amsart}
\usepackage{amsmath,amssymb,amsthm,graphicx}
\usepackage{verbatim}
\newenvironment{reference}{\comment}{\endcomment}
%\newenvironment{reference}{}{}
\newenvironment{slogan}{\comment}{\endcomment}
\newenvironment{history}{\comment}{\endcomment}

% For commutative diagrams we use Xy-pic
\usepackage[all]{xy}

% We use 2cell for 2-commutative diagrams.
\xyoption{2cell}
% Theorem environments.
%
\theoremstyle{plain}
\newtheorem{theorem}[subsection]{Theorem}
\newtheorem{proposition}[subsection]{Proposition}
\newtheorem{lemma}[subsection]{Lemma}

\theoremstyle{definition}
\newtheorem{definition}[subsection]{Definition}
\newtheorem{example}[subsection]{Example}
\newtheorem{exercise}[subsection]{Exercise}
\newtheorem{situation}[subsection]{Situation}

\theoremstyle{remark}
\newtheorem{remark}[subsection]{Remark}
\newtheorem{remarks}[subsection]{Remarks}

\numberwithin{equation}{subsection}

% Macros
%
\def\lim{\mathop{\mathrm{lim}}\nolimits}
\def\colim{\mathop{\mathrm{colim}}\nolimits}
\def\Spec{\mathop{\mathrm{Spec}}}
\def\Hom{\mathop{\mathrm{Hom}}\nolimits}
\def\Ext{\mathop{\mathrm{Ext}}\nolimits}
\def\SheafHom{\mathop{\mathcal{H}\!\mathit{om}}\nolimits}
\def\SheafExt{\mathop{\mathcal{E}\!\mathit{xt}}\nolimits}
\def\Sch{\mathit{Sch}}
\def\Mor{\mathop{\mathrm{Mor}}\nolimits}
\def\Ob{\mathop{\mathrm{Ob}}\nolimits}
\def\Sh{\mathop{\mathit{Sh}}\nolimits}
\def\NL{\mathop{N\!L}\nolimits}
\def\CH{\mathop{\mathrm{CH}}\nolimits}
\def\proetale{{pro\text{-}\acute{e}tale}}
\def\etale{{\acute{e}tale}}
\def\QCoh{\mathit{QCoh}}
\def\Ker{\mathop{\mathrm{Ker}}}
\def\Im{\mathop{\mathrm{Im}}}
\def\Coker{\mathop{\mathrm{Coker}}}
\def\Coim{\mathop{\mathrm{Coim}}}

% Boxtimes
%
\DeclareMathSymbol{\boxtimes}{\mathbin}{AMSa}{"02}

%
% Macros for moduli stacks/spaces
%
\def\QCohstack{\mathcal{QC}\!\mathit{oh}}
\def\Cohstack{\mathcal{C}\!\mathit{oh}}
\def\Spacesstack{\mathcal{S}\!\mathit{paces}}
\def\Quotfunctor{\mathrm{Quot}}
\def\Hilbfunctor{\mathrm{Hilb}}
\def\Curvesstack{\mathcal{C}\!\mathit{urves}}
\def\Polarizedstack{\mathcal{P}\!\mathit{olarized}}
\def\Complexesstack{\mathcal{C}\!\mathit{omplexes}}
% \Pic is the operator that assigns to X its picard group, usage \Pic(X)
% \Picardstack_{X/B} denotes the Picard stack of X over B
% \Picardfunctor_{X/B} denotes the Picard functor of X over B
\def\Pic{\mathop{\mathrm{Pic}}\nolimits}
\def\Picardstack{\mathcal{P}\!\mathit{ic}}
\def\Picardfunctor{\mathrm{Pic}}
\def\Deformationcategory{\mathcal{D}\!\mathit{ef}}
\usepackage{cleveref}
\crefname{originallemma}{original lemma}{original lemmas}
\Crefname{originallemma}{Original lemma}{Original lemmas}
\makeatletter
\newlabel{lemma-etale-makes-quasi-finite-finite-one-prime}{{00UJ}{source}{}{sourcepointer-1}{}}
\newlabel{lemma-etale-makes-quasi-finite-finite-one-prime@cref}{{[originallemma][0][]00UJ}{[1][0][]source}{}{}{}}
\makeatother
\title{After an etale base change a finite type algebra splits into a finite part and a part quasi-finite nowhere over the point}
\author{The Stacks Project Authors}
\date{stacks-project:algebra.tex:local-snapshot-20261012; GNU FDL 1.2 or later}
\begin{document}
\maketitle
\noindent\textit{Editorial note.} All delivered mathematical text is the author's own, copied from the pinned Stacks Project chapter file algebra.tex; no statement or proof step was rewritten, reordered or removed.\par\medskip
\noindent\textit{Editorial note.} Preamble provenance, as required for this source: the frozen snapshot stages the chapter's own preamble file beside it (preamble.tex, 217 lines, 8116 bytes, the same shared Stacks-wide preamble the published chow and etale-cohomology cards use); the first line of the pinned chapter file is \textbackslash{}input\{preamble\}, so no preamble text lives inside the chapter itself. The card takes that file's lines 8-18 (the verbatim package and the three \textbackslash{}comment-based environments reference/slogan/history at lines 9-12, section 13.2, together with the Xy-pic setup) and its lines 154 to the runtime end of the file (the native theorem declarations and mathematical macros). Not one character was changed, no package was added, no macro was defined, and the excerpt serves only as preamble, never as a span coordinate.\par\medskip
\noindent\textit{Editorial note.} The same-chapter pointers are located at their \textbackslash{}label lines inside the pinned chapter file; the cross-chapter pointers are located in the project's own tag list staged in the snapshot (sources/stacks-chow/tags.txt), at the verbatim lines carrying their canonical tags.\par\medskip
\noindent\textit{Editorial note.} The counted proof refers to same-chapter statements by label; they are recorded as navigation stubs pointed at their \textbackslash{}label lines, and the arguments using them stay inside the counted span.\par\medskip
\noindent\textit{Editorial note.} The counted spans contain no \textbackslash{}cite, so the card carries no bibliography.\par\medskip
\noindent\textit{Allowed standard prerequisites (editorial).}\begin{itemize}
\item The fibre ring is Noetherian with finitely many isolated closed points. Used at the start of the counted proof (lines 40603-40606). Separates the isolated points of the fibre, which are the pieces that can be made finite.
\item The structure of a finite algebra over a field at an isolated point. Used in the counted proof (the splitting of the fibre). Produces the product decomposition of the fibre at the isolated points.
\end{itemize}\medskip
\ifdefined\hypertarget\hypertarget{sourcepointer-1}{}\fi
\noindent\textit{Source reference.} Original lemma 00UJ: sources/stacks-algebra/pinned-original-source.tex, source lines 40445--40445. Navigation pointer for \textbackslash{}ref\{lemma-etale-makes-quasi-finite-finite-one-prime\} used inside the counted span (line 40610); the argument that consumes it is retained in the span.\par\medskip
\begin{lemma}
\label{lemma-etale-makes-quasi-finite-finite}
Let $R \to S$ be a ring map.
Let $\mathfrak p \subset R$ be a prime.
Assume $R \to S$ finite type.
Then there exists
\begin{enumerate}
\item an \'etale ring map $R \to R'$,
\item a prime $\mathfrak p' \subset R'$ lying over $\mathfrak p$,
\item a product decomposition
$$
R' \otimes_R S = A_1 \times \ldots \times A_n \times B
$$
\end{enumerate}
with the following properties
\begin{enumerate}
\item we have $\kappa(\mathfrak p) = \kappa(\mathfrak p')$,
\item each $A_i$ is finite over $R'$,
\item each $A_i$ has exactly one prime $\mathfrak r_i$ lying over
$\mathfrak p'$, and
\item $R' \to B$ not quasi-finite at any prime lying over $\mathfrak p'$.
\end{enumerate}
\end{lemma}

\begin{proof}
Denote $F = S \otimes_R \kappa(\mathfrak p)$ the fibre ring of $S/R$
at the prime $\mathfrak p$. As $F$ is of finite type over $\kappa(\mathfrak p)$
it is Noetherian and hence $\Spec(F)$ has finitely many isolated closed
points. If there are no isolated closed points,
i.e., no primes $\mathfrak q$ of $S$ over $\mathfrak p$ such that
$S/R$ is quasi-finite at $\mathfrak q$, then the lemma holds.
If there exists at least one such prime $\mathfrak q$, then
we may apply Lemma \ref{lemma-etale-makes-quasi-finite-finite-one-prime}.
This gives a diagram
$$
\xymatrix{
S \ar[r] & R'\otimes_R S \ar@{=}[r] & A_1 \times B' \\
R \ar[r] \ar[u] & R' \ar[u] \ar[ru]
}
$$
as in said lemma. Since the residue fields at $\mathfrak p$ and $\mathfrak p'$
are the same, the fibre rings of $S/R$ and $(A_1 \times B')/R'$
are the same. Hence, by induction on the number of isolated closed points
of the fibre we may assume that the lemma holds for
$R' \to B'$ and $\mathfrak p'$. Thus we get an \'etale ring
map $R' \to R''$, a prime $\mathfrak p'' \subset R''$ and
a decomposition
$$
R'' \otimes_{R'} B' = A_2 \times \ldots \times A_n \times B
$$
We omit the verification that the ring map $R \to R''$, the
prime $\mathfrak p''$ and the resulting decomposition
$$
R'' \otimes_R S = (R'' \otimes_{R'} A_1) \times
A_2 \times \ldots \times A_n \times B
$$
is a solution to the problem posed in the lemma.
\end{proof}

\end{document}
